\begin{table}[ht]
\centering
\small
\caption{Table A1. Main results, H1-H4 (fractional logit, appendix)}
\begin{tabular}{lrrrrrrrrrrrrr}
\toprule
Hypothesis & Exp. sign & N & M2 estimate & Scale & SE & 95\% CI & p (1-sided) & Holm p (1-sided) & Gold n & M2 on gold universe (logit) & PPI++ (logit) & PPI++ 95\% CI & PPI++ p (2-sided, Holm) \\
\midrule
H1 Opposition vs government askers & - & 144,349 & -0.003 & AME & 0.003 & {[}-0.009, +0.004] & 0.220 & 0.660 & 400 & -0.021 & +0.046 & {[}-0.173, +0.265] & 1.000 \\
H2 Coalition-dependent vs majority (LS18 vs LS17) & + & 94,981 & -0.031 & AME & 0.006 & {[}-0.043, -0.019] & 1.000 & 1.000 & 180 & -0.190 & -0.327 & {[}-0.648, -0.005] & 0.187 \\
H3 Security vs other ministries & - & 144,349 & -0.209 & logit & 0.033 & {[}-0.275, -0.143] & <0.001 & <0.001 & 400 & -0.211 & -0.229 & {[}-0.559, +0.100] & 0.510 \\
H4 Opposition x coalition-dependent (LS18 vs LS17) & + & 94,981 & -0.129 & logit & 0.029 & {[}-0.189, -0.068] & 1.000 & 1.000 & 180 & -0.138 & -0.134 & {[}-0.606, +0.337] & 1.000 \\
\bottomrule
\end{tabular}
\par\footnotesize Fractional logit with the same specifications and clustering as Table 3. M2 estimates for H1 and H2 are average marginal effects (score units); H3 and H4 are logit-scale coefficients (H4 is the interaction). PPI++ is estimated on the logit scale throughout; the logit-scale M2 estimate on the gold universe is shown for comparison. One-sided p and Holm as in Table 3.
\end{table}
